\documentclass[11pt]{article}
\usepackage[margin=0.9in]{geometry}
\usepackage{graphicx,amsmath,amssymb,booktabs,hyperref,xcolor}
\graphicspath{{/tmp/cheby_h0/figs/orders_report/}}
\hypersetup{colorlinks=true,linkcolor=blue,urlcolor=blue}
\setlength{\parskip}{6pt}
\title{Lin-CDF Richardson order sweep:\\
       R2 vs R3 vs R4 vs R5 across 100 $\gamma$ values}
\author{Fixed-Point Factory --- order vs accuracy vs robustness}
\date{June 2026}

\begin{document}
\maketitle

\section*{Headline}

Same architecture (Lin-CDF + Richardson + tabulated $\mu$), four orders
of Richardson extrapolation on the lookup table:

\begin{center}
\begin{tabular}{cccccc}
\toprule
order & $h$ values & $\sum |w_i|$ & machine $\varepsilon$ & slope @ asymptote & def$_{1{-}1}$ @ asymptote \\
\midrule
R2 & $\{0.5, 0.3\}$ & 2.12 & \textbf{100/100} & 0.508 & 0.001 \\
R3 & $\{0.5, 0.4, 0.3\}$ & 8.14 & 91/100 & 0.652 & 0.014 \\
R4 & $\{0.5, 0.4, 0.3, 0.2\}$ & 7.10 & \textbf{99/100} & 0.739 & 0.0001 \\
R5 & $\{0.6, 0.5, 0.4, 0.3, 0.2\}$ & 9.87 & 90/100 & 0.761 & \textbf{0.21}* \\
\bottomrule
\end{tabular}
\end{center}

(* R5 deficit at high $\gamma$ jumps to $0.21$ --- numerical
amplification from large weights, not a real economic effect.)

\paragraph{Reading.}
\begin{itemize}
  \item \textbf{R2}: bullet-proof, $100/100$ machine $\varepsilon$ at
        machine $\varepsilon$, but the deepest h-bias remains ($O(h^4)$).
  \item \textbf{R3}: slope jumps significantly (0.51 $\to$ 0.65) — the
        extra Richardson point pulls the FP toward a more
        T-revealing limit. Loses 9 gammas to numerical noise.
  \item \textbf{R4}: best practical sweet spot. 99/100 converged,
        slope rises further to 0.74, deficit collapses to $\sim 0$.
  \item \textbf{R5}: weights amplify too much, instability at high $\gamma$.
        Slope only marginally above R4 (0.76 vs 0.74) but deficit
        explodes to 0.21 at large $\gamma$.
\end{itemize}

\section{Slope $\alpha^*$ across orders}

\begin{center}
\includegraphics[width=\linewidth]{01_slope_orders.png}
\end{center}

Higher Richardson order pulls $\alpha^*$ up — each extra cancellation
removes more kernel bias from the lookup table. R2 plateaus around
$0.5$; R3 around $0.65$; R4/R5 around $0.74$--$0.76$. R4 and R5 essentially
agree on the slope, suggesting the Richardson series has converged.

\section{One-to-one deficit across orders}

\begin{center}
\includegraphics[width=\linewidth]{02_deficit_orders.png}
\end{center}

R2, R3, R4 all show \emph{decaying} deficit with $\gamma$ — the FP
approaches a function of $T$ at high $\gamma$. R5 has good behaviour
at low $\gamma$ but \emph{explodes} at high $\gamma$ (numerical
instability from amplification).

\section{Convergence floor}

\begin{center}
\includegraphics[width=\linewidth]{03_floor_orders.png}
\end{center}

R2 achieves 100/100 at machine $\varepsilon$; R4 nails 99/100. R3 and R5
have more failures (mostly in mid-$\gamma$).

\section{Why higher orders are unstable: Richardson weight magnitudes}

\begin{center}
\includegraphics[width=\linewidth]{04_weights.png}
\end{center}

\paragraph{The trade-off.} Each additional Richardson point cancels an
extra factor of $h^2$ in the bias, but the Vandermonde-derived weights
grow with order. The bilinear-interpolation error in each individual
$\mu^{h_i}$ table is amplified by $\sum |w_i|$ in the final extrapolation.

For Lin-CDF with bilinear interp:
\begin{itemize}
  \item R2: $\sum|w|=2.1$ — safe, but only $O(h^4)$ residual.
  \item R3: $\sum|w|=8.1$ — already amplifying. 9 failures.
  \item R4: $\sum|w|=7.1$ (h-spacing matters more than count). Best
        balance — only 1 failure, $O(h^8)$ residual.
  \item R5: $\sum|w|=9.9$ — too aggressive. 10 failures and instability
        at large $\gamma$.
\end{itemize}

For Chebyshev with polynomial interp (much smaller per-table error),
higher Richardson orders are tolerated.

\section{Recommendation}

\paragraph{For paper-level economics:} Use \textbf{R2} (most robust,
100\% machine $\varepsilon$, modest $h$-bias residual) OR \textbf{R4}
(best deeper-canceling order before instability hits).

\paragraph{The Richardson series HAS converged:} R4 and R5 agree on
slope ($0.74$ vs $0.76$) to within $\sim 3\%$, and R4 has a stable
near-zero deficit. So R4 is the recommended choice for high-fidelity
strict-$h{=}0$ on Lin-CDF.

\paragraph{Compare to Cheb:} Cheb tolerates 10-point Richardson because
polynomial interp has near-machine-precision error per $\mu^{h_i}$.
Lin-CDF's bilinear interp limits us to R4 at most. The trade-off is the
opposite of the speed comparison: Lin-CDF needs more careful Richardson
order tuning, but each individual operator call is faster.

\end{document}
